\subsection{问题四：考虑径向收缩的移动域模型}
\subsubsection{收缩几何与材料运动假设}
第四问外半径不再固定，而由附件2给出。数据以1800 s为间隔覆盖72 h，半径单位由厘米转换为米后线性插值，记为$R(t)$。药材域变为
\begin{equation}
\Omega(t)=\{0\le r\le R(t)\},\qquad
A_s(t)=2\pi R(t)L,\qquad V(t)=\pi R(t)^2L .
\end{equation}
测得半径只能确定外表面的运动，不能唯一确定内部骨架每个位置的速度。结合材料整体连续收缩的宏观描述\cite{mayor}，本文采用均匀径向运动
\begin{equation}
v(r,t)=\frac{\dot R(t)}{R(t)}\,r,\qquad
v(0,t)=0,\quad v(R(t),t)=\dot R(t).
\label{eq:velocity}
\end{equation}
这一假设使材料点的相对径向位置保持不变，适合通过归一化坐标描述移动域。由于$R(t)$分段线性，$\dot R$在采样折点处可跳变；转换后的方程只直接需要连续的$R(t)$，不要求对噪声半径数据额外数值求导。

\begin{figure}[H]\centering
\begin{tikzpicture}[scale=.95]
\draw[thick,ink] (0,0) circle (1.65);
\draw[teal] (0,0) circle (.825);
\fill[teal](.825,0) circle (2pt);
\draw[-{Stealth},ink] (0,0)--(1.65,0);
\node at (0,-2.0) {初始域：$r_0=xR_0$};
\draw[arr](2.0,.2)--(4.0,.2);
\node[above,align=center] at (3,0.3) {均匀径向收缩\\$x$不变};
\draw[thick,ink] (5.5,0) circle (1.1);
\draw[teal] (5.5,0) circle (.55);
\fill[teal](6.05,0) circle (2pt);
\draw[-{Stealth},ink](5.5,0)--(6.6,0);
\node at (5.5,-2.0) {当前域：$r(t)=xR(t)$};
\node[above] at (.65,.03) {$x=0.5$};
\node[above] at (5.9,.03) {$x=0.5$};
\end{tikzpicture}
\caption{材料坐标与实际距离的对应关系示意，圆的尺寸仅用于说明映射}\label{fig:shrinkmap}
\end{figure}

固定实际位置和固定材料点是两种不同观测。固定$r=1$ cm时，材料可能因收缩从该处移走；固定$x=0.5$时则始终跟随同一骨架位置，其实际距离随$R(t)$改变。第四问的计算在材料坐标中完成，题目结果必须再转换回实际距离。

收缩使扩散路径缩短，并改变表面积与体积比，但第四问同时更换全部经验物性。因此第三、四问时长差是两个模型方案的综合差异。若要单独识别收缩贡献，需要固定物性作额外对照；本次依照要求不新增计算。

\nextpage
\subsubsection{骨架、水分连续性与坐标变换}
在固定空间坐标中，单位体积的干物质量满足
\begin{equation}
(\rho_d)_t+\frac1r\pd_r(r\rho_dv)=0 .
\label{eq:skeleton}
\end{equation}
水分同时随骨架运动和相对骨架扩散，总连续性方程为
\begin{equation}
(\rho_dC)_t+\frac1r\pd_r[r(\rho_dCv+j)]=0,
\qquad j=-\rho_dD C_r .
\end{equation}
展开后减去$C$乘以式\eqref{eq:skeleton}，得到
\begin{equation}
\rho_d(C_t+vC_r)=\frac1r\pd_r(r\rho_dD C_r).
\label{eq:materialwater}
\end{equation}
因此，收缩不需要向含水率方程额外添加任意“体积浓缩源”。对于同一干物质量，$C$表示水与干物质之比，随骨架的共同压缩已在连续性中一致处理。若$\rho_d$空间均匀，结合式\eqref{eq:velocity}有$\rho_d(t)R(t)^2=\text{常数}$，可以从式\eqref{eq:materialwater}中约去$\rho_d$。

引入$x=r/R(t)$及$\widehat f(x,t)=f(R(t)x,t)$。由链式法则
\begin{equation}
f_r=\frac{\widehat f_x}{R},\qquad
\widehat f_t=f_t+x\dot R f_r ,
\end{equation}
从而
\begin{equation}
f_t+vf_r=
\widehat f_t+\frac{v(Rx,t)-x\dot R}{R}\widehat f_x
=\widehat f_t .
\end{equation}
最后一个等号只在均匀径向收缩假设下成立。扩散算子转换为
\begin{equation}
\frac1r\pd_r(rDf_r)=
\frac{1}{R(t)^2x}\pd_x(xD\widehat f_x).
\label{eq:transform}
\end{equation}
于是问题域固定为$0\le x\le1$，材料运动已经体现在坐标映射中\cite{huang}。这一过程明确区分物质导数与固定空间偏导，避免把网格移动与真实水分扩散混为一项。

温度采用随骨架运动的有效显热近似，以$\rho c_p(T_t+vT_r)$作为升温项；仍忽略潜热、焓迁移和机械变形功。它与水分精确的参考干质量连续性各有物理约定，不能据此推断题给经验密度同时满足所有湿质量约束。

\nextpage
\subsubsection{变域耦合模型与数值实现}
题目附录4的物性为
\begin{equation}
\begin{aligned}
\rho(C)&=760+90C,&c_p(C)&=1850+2150\frac{C}{1+C},\\
k(C)&=0.12+0.20\frac{C}{1+C},&
D(C,T_K)&=4.2\times10^{-4}
\exp\left(-\frac{0.30}{C}-\frac{3850}{T_K}\right).
\end{aligned}
\label{eq:props4}
\end{equation}
结合式\eqref{eq:transform}，实际求解模型为
\begin{equation}
\begin{cases}
\rho(\widehat C)c_p(\widehat C)\widehat T_t
=\dfrac{1}{R(t)^2x}\pd_x[xk(\widehat C)\widehat T_x],\\[5pt]
\widehat C_t
=\dfrac{1}{R(t)^2x}\pd_x[xD(\widehat C,\widehat T_K)\widehat C_x].
\end{cases}
\label{eq:q4}
\end{equation}
初始状态为$\widehat T(x,0)=28,\widehat C(x,0)=2.55$，中心对称条件为$\widehat T_x(0,t)=\widehat C_x(0,t)=0$。表面条件必须同时换算空间导数：
\begin{equation}
-\frac{k_s}{R(t)}\widehat T_x(1,t)
=h(\widehat T_s-T_a),\qquad
-\frac{D_s}{R(t)}\widehat C_x(1,t)
=h_m(\widehat C_s-C_e).
\label{eq:q4bc}
\end{equation}
漏掉任一$1/R(t)$因子，都会改变表面传递强度。离散时固定$x_i$，在每次右端评估中读取当前$R(t)$，内部通量按$R^{-2}$缩放，表面通量按$R^{-1}$缩放；共享通量和控制体更新形式与问题一一致。

第四问仍使用$N=2560,\beta=5$和相对容差$2\times10^{-10}$，并沿用4 h前后分段环境。事件函数为$\max_i\widehat C_i(t)-0.15$，终止后在实际空间位置$r_q$取值：
\begin{equation}
C(r_q,t)=\widehat C\!\left(\frac{r_q}{R(t)},t\right),
\qquad 0\le r_q\le R(t).
\end{equation}
空间插值使用当前时刻的比例坐标，不能把初始$r_q/R_0$始终当作同一个实际距离。如果$r_q>R(t)$，该点不属于当前药材内部，结果留空；另设“药材表面”列，恒取$x=1$。

收缩后$\rho_d$虽然随时间增加，但对归一化控制体而言，参考干质量$\rho_d\,2\pi LR(t)^2\omega_i$保持不变。故$\overline C=2\sum_i\omega_i C_i$仍可解释为干质量加权平均，离散水分守恒检查保持成立。

\nextpage
\subsubsection{收缩后的干燥时长与实际位置结果}
\begin{table}[H]\centering
\caption{问题4：实际位置的干基含水率（kg/kg）与当前半径}\label{tab:q4}
\begin{tabular}{cccccc}\toprule
时间/h & 0 cm & 0.5 cm & 1 cm & 表面 & $R$/cm\\ \midrule
6.0000 & 1.7196 & 1.5378 & 1.0226 & 0.4207 & 1.3740\\
12.0000 & 0.7377 & 0.6546 & 0.4083 & 0.1669 & 1.2480\\
18.0000 & 0.4088 & 0.3687 & 0.2397 & 0.0892 & 1.2140\\
24.0000 & 0.2851 & 0.2611 & 0.1790 & 0.0673 & 1.2040\\
30.0000 & 0.2263 & 0.2094 & 0.1493 & 0.0595 & 1.2010\\
36.0000 & 0.1928 & 0.1797 & 0.1317 & 0.0560 & 1.2000\\
42.0000 & 0.1712 & 0.1604 & 0.1200 & 0.0541 & 1.2000\\
48.0000 & 0.1561 & 0.1468 & 0.1116 & 0.0530 & 1.2000\\
51.0920 & 0.1500 & 0.1413 & 0.1081 & 0.0526 & 1.2000\\
\bottomrule\end{tabular}\end{table}


主计算得到
\begin{equation}
t_*^{(4)}=183931.2084622\ \mathrm s
=51.0920023506\ \mathrm h ,
\end{equation}
按四位小数报告为\textbf{51.0920 h}。临界时最大含水率位于中心，表面为0.0526 kg/kg，当前半径为1.2000 cm。若按整秒满足严格低于要求，可采用183932 s；若以四位小数小时并向上取整，则为51.0921 h。正文表中的51.0920是临界时间的常规舍入。

表\ref{tab:q4}各展示时刻的半径均小于1.5 cm，所以只报告0、0.5、1 cm三个固定位置和当前表面，符合题目允许按实际尺寸展开位置列的含义。临界时半径虽然约为初始的60\%，但不能据此直接按$R^2$将时长缩放为原来的36\%；因为整个过程中半径在变化，扩散系数也受到温湿状态及物性方案影响。

第三、四问的临界时长差为
\begin{equation}
\Delta t_*=t_*^{(3)}-t_*^{(4)}
\approx6.3820\ \mathrm h,\qquad
\frac{\Delta t_*}{t_*^{(3)}}\approx11.10\%.
\end{equation}
在6 h时第四问中心含水率反而高于第三问，而临界终点更早，说明不能把整体比较压缩成“收缩必然在所有时刻都使内部更干”。正确结论是：在题给新物性和均匀径向收缩共同作用下，第四问最终临界时间较短。

\texttt{result4.xlsx}按60 s、0.1 cm保存实际位置水分，越出当前域的位置留空，并附表面含水率及当前半径。空值表示“没有材料”，零值则表示“有材料但水分为零”，二者不可替换。
